% s-084-p084-m003; original PDF 84, 85
\Aphorism{s-084-p084-m003}{80}
Es giebt eine peinliche Art von „Idealisten“, die sich selber in Versuchung führen, um sich als Unterliegende zu ironisiren, die beständig beim Sumpfe wohnen, um sich beständig über schlechte Luft zu beschweren, die mit Eigensinn das ihnen Widerwärtige aufsuchen, nur um mit dem Genusse des Schimpfens sich und Anderen die Ohren vollzugellen. Unsere Gegenwart bringt diese Pilze und Sumpfblumen in Fülle hervor; mit Vorliebe halten sie sich in der Nähe Gottes, der Gesellschaft, des Staates auf. Manche von ihnen waren ursprünglich edel und sind nur durch Noth und Abhängigkeit verdorben, andere haben noch etwas Zukunft; die meisten sind ohne Hoffnung wie Tradition, kleinliche giftige Kritiker ihrer localen und zeitlichen Nachbarschaft, unmögliche Menschen, die es schlecht haben \Sperr{müssen}, um das Leben zu fristen. Gegen sie hat Zarathustra sein edelstes Wort gesprochen: wo man nicht mehr lieben kann, da soll man \Sperr{vorübergehn}! Der erste dieser Idealisten ist übrigens Gott selbst, jener Gott, der die Sünde hasste und dennoch die Sünde schuf. Warum ist er nicht \Sperr{vorübergegangen}?
\par

